
\chapter{Metodologi Penelitian}

%=========================================================%
%             TULIS METODE PENELITIAN DI SINI             %
% Jangan Lupa untuk Menghapus Contoh Tulisan di Bawah Ini %
%          (Termasuk Panduan di Setiap Subbabnya)         %
% Contoh Teks Dapat Di-comment dengan Tanda "%" di Depan  %
%   jika Ingin Mempertahankan Panduan di Dalam File Ini   %
%=========================================================%

% --- Pendekatan Penelitian ---
%
%     Dapat digabung dengan subbab berikutnya menjadi
%     "Pendekatan dan Jenis Penelitian" agar lebih
%     efisien jika mau. Jangan lupa menghapus atau
%     memberi comment pada \section{Jenis Penelitian}
%     bila memutuskan untuk menggabung keduanya.
%
\section{Pendekatan Penelitian}

Uraikan pendekatan penelitian yang Anda gunakan, apakah paradigma penelitian Anda condong ke interpretasi objektif (kuantitatif/positivistik) atau subjektif (kualitatif/konstruktif).



% --- Jenis Penelitian ---
%
%     Dapat dihapus bila ingin digabung ke subbab sebelumnya.
%
\section{Jenis Penelitian}

Selanjutnya, jabarkan jenis penelitiannya, seperti eksplanatif (menjelaskan hubungan sebab-akibat), deskriptif (memaparkan fenomena secara mendalam), eksploratif, atau evaluatif. Berikan alasan metodologis yang kuat mengapa pendekatan dan jenis tersebut paling tepat untuk menjawab masalah penelitian Anda.



% --- Metode Penelitian ---
\section{Metode Penelitian}

Jelaskan secara spesifik desain atau metode operasional yang dipilih untuk menjalankan penelitian ini. Contoh metode yang bisa digunakan antara lain: survei, eksperimen, analisis isi (\textit{content analysis}), fenomenologi, studi kasus, atau metode rekayasa sistem yang relevan. Berikan argumen yang rasional mengapa metode ini dipilih dan bagaimana metode ini mampu memandu Anda dalam mengumpulkan data secara sistematis.



% --- Lokasi dan Waktu Penelitian ---
\section{Lokasi dan Waktu Penelitian}

Deskripsikan secara jelas di mana penelitian ini dilaksanakan (lokasi fisik, instansi, perusahaan, atau batas geografis/digital tertentu). Selain lokasi, jelaskan linimasa (waktu) pelaksanaan penelitian mulai dari tahap persiapan, pengumpulan data, analisis, hingga penyusunan laporan. Boleh juga untuk menyajikan linimasa ini dalam bentuk tabel jadwal kegiatan (\textit{Gantt chart}) agar terlihat bagus.



% --- Teknik Pengumpulan Data ---
%
%     Judul dapat diganti menjadi "Populasi, Sampel, dan Teknik
%     Pengumpulan" jika mau.
%
%     Bila berkaitan dengan teknik komputer, Teknik Pengumpulan
%     Data dapat disesuaikan menjadi "Teknik Pengumpulan
%     Kebutuhan"
%
\section{Teknik Pengumpulan Data}

Subbab ini dibagi menjadi beberapa perintilan untuk menjelaskan subjek dan data Anda:

\begin{enumerate}
    \item \textbf{Populasi dan Sampel:} Deskripsikan karakteristik populasi target, batasan sampel, serta siapa respondennya. Jelaskan pula proses penentuan \textit{sampling} (misalnya \textit{probability} atau \textit{non-probability sampling}) secara gamblang.
    
    \item \textbf{Teknik Pengumpulan Data:} Paparkan instrumen apa saja yang digunakan untuk menjaring informasi (seperti kuesioner, wawancara mendalam, \textit{Focus Group Discussion} (FGD), studi pustaka, atau observasi langsung).
    
    \item \textbf{Pengembangan Instrumen:} Jelaskan bagaimana Anda menyusun instrumen tersebut (misalnya pembuatan kisi-kisi pertanyaan wawancara atau pengujian validitas dan reliabilitas butir kuesioner).
\end{enumerate}



% --- Teknik Analisis Data ---
\section{Teknik Analisis Data}

Uraikan secara runtut cara Anda mengolah, menyajikan, dan menginterpretasikan data yang telah dikumpulkan. Teknik analisis data sangat bergantung pada metode penelitian Anda: bisa berupa analisis statistik dengan bantuan perangkat lunak (seperti SPSS, R, SmartPLS), atau analisis teks kualitatif (seperti analisis semiotika, \textit{framing}, analisis tematik, atau analisis konten). Pastikan langkah-langkah analisis ditulis secara logis dari reduksi data hingga penarikan kesimpulan.